% Options for packages loaded elsewhere
\PassOptionsToPackage{unicode}{hyperref}
\PassOptionsToPackage{hyphens}{url}
\PassOptionsToPackage{dvipsnames,svgnames,x11names}{xcolor}
%
\documentclass[
  letterpaper,
  DIV=11,
  numbers=noendperiod]{scrartcl}

\usepackage{amsmath,amssymb}
\usepackage{iftex}
\ifPDFTeX
  \usepackage[T1]{fontenc}
  \usepackage[utf8]{inputenc}
  \usepackage{textcomp} % provide euro and other symbols
\else % if luatex or xetex
  \usepackage{unicode-math}
  \defaultfontfeatures{Scale=MatchLowercase}
  \defaultfontfeatures[\rmfamily]{Ligatures=TeX,Scale=1}
\fi
\usepackage{lmodern}
\ifPDFTeX\else  
    % xetex/luatex font selection
\fi
% Use upquote if available, for straight quotes in verbatim environments
\IfFileExists{upquote.sty}{\usepackage{upquote}}{}
\IfFileExists{microtype.sty}{% use microtype if available
  \usepackage[]{microtype}
  \UseMicrotypeSet[protrusion]{basicmath} % disable protrusion for tt fonts
}{}
\makeatletter
\@ifundefined{KOMAClassName}{% if non-KOMA class
  \IfFileExists{parskip.sty}{%
    \usepackage{parskip}
  }{% else
    \setlength{\parindent}{0pt}
    \setlength{\parskip}{6pt plus 2pt minus 1pt}}
}{% if KOMA class
  \KOMAoptions{parskip=half}}
\makeatother
\usepackage{xcolor}
\setlength{\emergencystretch}{3em} % prevent overfull lines
\setcounter{secnumdepth}{5}
% Make \paragraph and \subparagraph free-standing
\makeatletter
\ifx\paragraph\undefined\else
  \let\oldparagraph\paragraph
  \renewcommand{\paragraph}{
    \@ifstar
      \xxxParagraphStar
      \xxxParagraphNoStar
  }
  \newcommand{\xxxParagraphStar}[1]{\oldparagraph*{#1}\mbox{}}
  \newcommand{\xxxParagraphNoStar}[1]{\oldparagraph{#1}\mbox{}}
\fi
\ifx\subparagraph\undefined\else
  \let\oldsubparagraph\subparagraph
  \renewcommand{\subparagraph}{
    \@ifstar
      \xxxSubParagraphStar
      \xxxSubParagraphNoStar
  }
  \newcommand{\xxxSubParagraphStar}[1]{\oldsubparagraph*{#1}\mbox{}}
  \newcommand{\xxxSubParagraphNoStar}[1]{\oldsubparagraph{#1}\mbox{}}
\fi
\makeatother


\providecommand{\tightlist}{%
  \setlength{\itemsep}{0pt}\setlength{\parskip}{0pt}}\usepackage{longtable,booktabs,array}
\usepackage{calc} % for calculating minipage widths
% Correct order of tables after \paragraph or \subparagraph
\usepackage{etoolbox}
\makeatletter
\patchcmd\longtable{\par}{\if@noskipsec\mbox{}\fi\par}{}{}
\makeatother
% Allow footnotes in longtable head/foot
\IfFileExists{footnotehyper.sty}{\usepackage{footnotehyper}}{\usepackage{footnote}}
\makesavenoteenv{longtable}
\usepackage{graphicx}
\makeatletter
\newsavebox\pandoc@box
\newcommand*\pandocbounded[1]{% scales image to fit in text height/width
  \sbox\pandoc@box{#1}%
  \Gscale@div\@tempa{\textheight}{\dimexpr\ht\pandoc@box+\dp\pandoc@box\relax}%
  \Gscale@div\@tempb{\linewidth}{\wd\pandoc@box}%
  \ifdim\@tempb\p@<\@tempa\p@\let\@tempa\@tempb\fi% select the smaller of both
  \ifdim\@tempa\p@<\p@\scalebox{\@tempa}{\usebox\pandoc@box}%
  \else\usebox{\pandoc@box}%
  \fi%
}
% Set default figure placement to htbp
\def\fps@figure{htbp}
\makeatother
% definitions for citeproc citations
\NewDocumentCommand\citeproctext{}{}
\NewDocumentCommand\citeproc{mm}{%
  \begingroup\def\citeproctext{#2}\cite{#1}\endgroup}
\makeatletter
 % allow citations to break across lines
 \let\@cite@ofmt\@firstofone
 % avoid brackets around text for \cite:
 \def\@biblabel#1{}
 \def\@cite#1#2{{#1\if@tempswa , #2\fi}}
\makeatother
\newlength{\cslhangindent}
\setlength{\cslhangindent}{1.5em}
\newlength{\csllabelwidth}
\setlength{\csllabelwidth}{3em}
\newenvironment{CSLReferences}[2] % #1 hanging-indent, #2 entry-spacing
 {\begin{list}{}{%
  \setlength{\itemindent}{0pt}
  \setlength{\leftmargin}{0pt}
  \setlength{\parsep}{0pt}
  % turn on hanging indent if param 1 is 1
  \ifodd #1
   \setlength{\leftmargin}{\cslhangindent}
   \setlength{\itemindent}{-1\cslhangindent}
  \fi
  % set entry spacing
  \setlength{\itemsep}{#2\baselineskip}}}
 {\end{list}}
\usepackage{calc}
\newcommand{\CSLBlock}[1]{\hfill\break\parbox[t]{\linewidth}{\strut\ignorespaces#1\strut}}
\newcommand{\CSLLeftMargin}[1]{\parbox[t]{\csllabelwidth}{\strut#1\strut}}
\newcommand{\CSLRightInline}[1]{\parbox[t]{\linewidth - \csllabelwidth}{\strut#1\strut}}
\newcommand{\CSLIndent}[1]{\hspace{\cslhangindent}#1}

\KOMAoption{captions}{tableheading}
\makeatletter
\@ifpackageloaded{caption}{}{\usepackage{caption}}
\AtBeginDocument{%
\ifdefined\contentsname
  \renewcommand*\contentsname{Table of contents}
\else
  \newcommand\contentsname{Table of contents}
\fi
\ifdefined\listfigurename
  \renewcommand*\listfigurename{List of Figures}
\else
  \newcommand\listfigurename{List of Figures}
\fi
\ifdefined\listtablename
  \renewcommand*\listtablename{List of Tables}
\else
  \newcommand\listtablename{List of Tables}
\fi
\ifdefined\figurename
  \renewcommand*\figurename{Figure}
\else
  \newcommand\figurename{Figure}
\fi
\ifdefined\tablename
  \renewcommand*\tablename{Table}
\else
  \newcommand\tablename{Table}
\fi
}
\@ifpackageloaded{float}{}{\usepackage{float}}
\floatstyle{ruled}
\@ifundefined{c@chapter}{\newfloat{codelisting}{h}{lop}}{\newfloat{codelisting}{h}{lop}[chapter]}
\floatname{codelisting}{Listing}
\newcommand*\listoflistings{\listof{codelisting}{List of Listings}}
\makeatother
\makeatletter
\makeatother
\makeatletter
\@ifpackageloaded{caption}{}{\usepackage{caption}}
\@ifpackageloaded{subcaption}{}{\usepackage{subcaption}}
\makeatother

\usepackage{bookmark}

\IfFileExists{xurl.sty}{\usepackage{xurl}}{} % add URL line breaks if available
\urlstyle{same} % disable monospaced font for URLs
\hypersetup{
  pdftitle={Smoking prevalence projection for Birmingham},
  pdfauthor={Chung Au-Yeung},
  colorlinks=true,
  linkcolor={blue},
  filecolor={Maroon},
  citecolor={Blue},
  urlcolor={Blue},
  pdfcreator={LaTeX via pandoc}}


\title{Smoking prevalence projection for Birmingham}
\author{Chung Au-Yeung}
\date{}

\begin{document}
\maketitle

\renewcommand*\contentsname{Table of contents}
{
\hypersetup{linkcolor=}
\setcounter{tocdepth}{3}
\tableofcontents
}

\section{Introduction}\label{introduction}

Birmingham Public Health has expanded community stop smoking services
through the Smokefree Generation grant, marking a major step forward in
the national journey towards becoming smokefree by 2030. The initiative
expands capacity for cessation services in communities where smoking
prevalence remains high and introduces targeted interventions for groups
disproportionately affected by tobacco use---particularly those
experiencing mental illness, substance use, or homelessness.

While these initiatives strengthen local cessation capacity, reducing
smoking prevalence across all communities requires sustained monitoring.
This report aims at developing a model to project the smoking prevalence
in Birmingham till 2047, allowing the local team to track progress
toward the 2030 target. The model will also allow scenario testing to
assess the potential impact of future policies or interventions.

\section{Method}\label{method}

\subsection{Data}\label{data}

This projection draws upon publicly available data from multiple
reputable sources. Population and migration estimates for Birmingham up
to 2047 were obtained from the Office for National Statistics (ONS)
2022-based subnational population projections {[}1{]}. Smoking
prevalence data were collated from several national datasets, including
Adult Smoking Habits in England (ONS) {[}2{]}, Smoking, Drinking and
Drug Use among Young People in England 2023 (NHS Digital) {[}3{]}, and
the Profile of 16--17-year-old Smokers from Action on Smoking and Health
(ASH) {[}4{]}.

{The projection has updated with the fololwing data. Population and
migration estimates for Birmingham up to 2047 were obtained from the
Office for National Statistics (ONS) using 2024-based subnational
population projections. {[}1{]}. Adult Smoking Habits in England (ONS)
{[}5{]}.}

In each simulation cycle, a new cohort of 13-year-olds is introduced as
non-smokers. As smoking prevalence at this age is negligible, therefore,
the model assumes no 13-year-olds smoke. The injection of new cohort of
13-year-olds for each simulation cycle is obtained from ONS subnational
projection {[}1{]}.

In each cycle, individuals can die from any states. The all-cause
mortality rates were derived from the national life tables 2021-2023
{[}6{]} and subsequently averaged across the three smoking status
groups: non-smokers, current smokers, and ex-smokers. This analysis uses
the equation from a prior study to calculate the death risk for
non-smoker {[}7{]}, as shown below:

\[ P_{d0} = \frac{P_{d\text{all}}}{P_{r\text{non}} + P_{r\text{cur}} \times R_{\text{cur}} + P_{r\text{ex}} \times R_{\text{ex}}} \]

\(P_{d\text{all}}\) is the all cause mortality rate while
\(P_{d\text{rnon}}\), \(P_{d\text{rcur}}\) and \(P_{d\text{rex}}\) are
prevalence of non-smoker, current smoker and ex-smoker respectively. The
mortality rate for current smoker are elevated by the relative risk
\(R_{d\text{cur}}\) and therefore, the rate is
\(P_{d0} \times R_{\text{cur}}\). Correspondingly, the mortality rate
for ex-smoker will be \(P_{d0} \times R_{\text{ex}}\). The relative
risks of death for current and ex-smokers were obtained from published
literature {[}8{]}. As relative risk vary according to duration since
quitting (e.g., short-term vs.~long-term cessation) and age at quitting,
there will have different combinations of relative risks. However, the
model used in this analysis does not account for age at cessation or
distinguish between short- and long-term quitters. Therefore, the
relative risk for ex-smokers was based on estimates corresponding to
long-term quitters. Furthermore, the model includes 10 sub-states for
ex-smokers, corresponding to the number of years since quitting. These
sub-states allow ex-smokers to experience a diminishing probability of
relapse over time, reflecting evidence that the likelihood of returning
to smoking decreases with the duration of abstinence.

The transition probabilities among states in the Markov model were
specific to age and gender. These probabilities were obtained an
estimation produced from cross-sectional survey by the University of
Sheffield {[}9{]}. The original probabilities were reported by Index of
Multiple Deprivation (IMD) quintile. In order to generate an overall
estimate representative of Birmingham, a weighted average was calculated
using the population distribution of smokers across IMD deciles within
the city. Since smoking overall has being following a downward trend in
recent years and this pattern is expected to persist, it is possible
that the smoking prevalence could plateau or even increase, similar to
trend being observed in Australia and US. Therefore, this analysis
follows the Department of Health \& Social Care (DHSC)'s method {[}10{]}
to fix the transition probabilities at 2022 estimate through out the
simulation {(The updated projection used probabilities at 2024
estimate).} With regards to DHSC's method again, a long-term quit
probability was applied to ex-smokers who had remained abstinent for 10
years, allowing them to transition to the non-smoker state. This implies
that a proportion of ex-smokers will share the same health risk as the
non-smokers.

In order to follow the DHSC's method as closely as possible, instigation
rates for 13- to 15-year-olds were calculated using a scaling approach.
This was carried out by referencing the instigation rates for males and
females from the US SimSmoke model {[}11{]}, applying the assumption
that the ratio between adjacent age groups in the US model also applies
to the Birmingham population.

\subsection{Model structure}\label{model-structure}

This analysis employed an open-cohort Markov model to simulate the
natural history of smoking behaviour for those who are above 13 years
old in Birmingham. The model includes annual cycles and a finite time
horizon of 25 years {(23 years for latest data)}. The model started at
time 0 using the number of population in 2022 {(2024 for latest data)}.
Smoking status was classified into non-smokers, current smokers and
ex-smokers. At the start of each annual cycle, the distribution of
smoking status among these entrants is determined in proportion to
observed smoking prevalence data for this age group. Simultaneously, the
model incorporates net migration by smoking status. For simplicity, it
is assumed that the smoking distribution of migrants mirrors that of the
UK population.

\begin{figure}[H]

{\centering \pandocbounded{\includegraphics[keepaspectratio]{images/statetransition.png}}

}

\caption{\textbf{Figure 1.} Markov state transition diagram for smoking
prevalence}

\end{figure}%

Figure 1 shows that at the beginning of each year t, the population is
divided into three main states: non-smokers, current smokers, and
ex-smokers.

\begin{itemize}
\item
  \textbf{Non-smokers:} At year \emph{t+1}, the number of non-smokers
  equals to (1) the number of people who were non-smokers in year
  \emph{t} and remained non-smokers, plus (2) the number of ex-smokers
  who reached 10 years since quitting in year \emph{t} and transitioned
  back to the non-smoker state, minus (3) the number of non-smokers in
  year \emph{t} who started smoking and moved to the smoker state in
  year \emph{t+1}, minus (4) the number of non-smokers in year \emph{t}
  who died, plus (5) the net number of non-smokers migrating into
  Birmingham (in-migrants minus out-migrants) during year \emph{t}.
\item
  \textbf{Current smokers:} At year \emph{t}+1, the number of current
  smokers equals (1) those who were smokers in year t and continued
  smoking, plus (2) non-smokers in year t who started smoking in year
  \emph{t}+1, plus (3) ex-smokers in year \emph{t} who relapsed and
  returned to smoking, minus (4) smokers in year \emph{t} who quit and
  moved to the ex-smoker state, minus (5) smokers in year \emph{t} who
  died, plus (6) the net number of smokers migrating into Birmingham
  during year \emph{t}.
\item
  \textbf{Ex-smokers:} At year \emph{t}+1, the number of ex-smokers
  equals (1) those who were ex-smokers in year \emph{t} and remained
  ex-smokers, plus (2) smokers in year \emph{t} who quit during the year
  and became new ex-smokers, plus (3) every former smoker who quit
  between year \emph{t--10} and \emph{t} and continues to remain
  smoke-free (i.e., those progressing through the ``time since quit''
  categories up to 10 years), minus (4) ex-smokers in any ``time since
  quit'' category who relapsed and returned to smoking, minus (5)
  ex-smokers in any ``time since quit'' category who died during year
  \emph{t}, minus (6) ex-smokers in year \emph{t} who reached 10 years
  since quitting and transitioned back to the non-smoker state, plus (7)
  the net number of ex-smokers migrating into the area during year
  \emph{t}.
\end{itemize}

At the end of each cycle, every person ages by one year (i.e., age
\emph{a} becomes \emph{a+1} in year \emph{t+1}), and in each new year
\emph{t+1}, a new cohort of 13-year-olds enters the model as
non-smokers. Every person is given a new set of age and gender specific
transition probabilities when cycle moves forward. For example, a
17-year-old male in year t will be assigned the transition probabilities
corresponding to 17-year-old males; when he becomes 18 in the next cycle
(\emph{t}+1), his probabilities are updated to those specific to
18-year-old males. This ensures that changes in smoking behaviour
reflect age- and gender-related differences in initiation, quitting,
relapse, and mortality rates over time.

\subsection{How scenario works?}\label{how-scenario-works}

To model each policy scenario, assumptions of reduction in instigation
rate were applied at 10\%, 30\%, 60\% and 90\%. It has to be noted that
the reduction in instigation is year-by-year, meaning that reduction
will get compounded over time, implying that the implementation of the
policy does not provide a one-off change but the effect becomes
progressively more effective. For example, under a 10\% annual reduction
scenario, the instigation rate in year \emph{t} would be multiplied by
0.9. In the following year (\emph{t+1}), the rate would be further
multiplied by 0.9², and so on, leading to a cumulative reduction effect
over time.

\section{Results}\label{results}

\subsection{Baseline result}\label{baseline-result}

Applying the baseline transition probabilities to the initial population
produces the baseline scenario, representing outcomes in the absence of
any policy intervention. The result is comparable to the DHSC's
published model, showing a similar trend and similar prevalence. The
first cycle of the model in 2023 shows the the smoking prevalence for
those among 14 to 30 year old is 11.35\%. Without any additional
policies being implemented, the smoking prevalence in Birmingham is
estimated to reach 11\% in 2047.

{From the latest data, the first cycle of the model in 2025 shows the
the smoking prevalence for those among 14 to 30 year old is 10.45\%.
Without any additional policies being implemented, the smoking
prevalence in Birmingham is estimated to reach 10.15\% in 2047.}

The trend is most likely similar to the published DHSC's model where
both models reached the long run equilibrium very quickly in the year of
2028 {(2029 for latest data)} but the change in the rate was
significantly slower for Birmingham. This difference can largely be
attributed to demographic and socio-economic factors. Birmingham has a
younger population structure, both in the 2022 {(2024 for latest data)}
baseline and in the subnational population projections to 2047 produced
by the ONS, compared with England overall. A younger age profile implies
a larger proportion of individuals in age groups with higher instigation
rates and lower long-term cessation, thereby moderating the overall
decline.

In addition, the quit and relapse probabilities in this model were
calculated as a weighted average by deprivation quintile. Given
Birmingham is one of the most deprived local authorities, the quit rate
is lower and relapse rate is higher relative to the national average.
Consequently, less individuals quit and more individuals transition back
from ex-smoker to current smoker states, resulting in a slower early
decline in prevalence before reaching equilibrium.

Moreover, this analysis employed an open-cohort markov model, meaning
that migration was incorporated, whereas the DHSC model was based on a
closed cohort. This might help explain why the decline in smoking
prevalence was slower in this analysis compared to the DHSC's model.

\pandocbounded{\includegraphics[keepaspectratio]{smokingprevalence_reading_document_files/figure-pdf/unnamed-chunk-1-1.png}}

The baseline result for adults (18+) is comparable to the model produced
by Cancer Research UK. The first cycle of the model in 2023 shows the
the smoking prevalence for adults is 13.06\%. Without any additional
policies being implemented, the smoking prevalence in Birmingham is
estimated to reach 6.4\% in 2047. While the trend is most likely similar
to the one from Cancer Research UK {[}12{]}, their model produced a
faster declining rate, the same reasons from the above can explain such
deviation.

{For the latest data, the first cycle of the model in 2025 shows the the
smoking prevalence for adults is 12.13\%. Without any additional
policies being implemented, the smoking prevalence in Birmingham is
estimated to reach 6.07\% in 2047.}

\pandocbounded{\includegraphics[keepaspectratio]{smokingprevalence_reading_document_files/figure-pdf/unnamed-chunk-2-1.png}}

\subsection{Scenario}\label{scenario}

This analysis followed the scenario where the Tobacoo and Vape bill
{[}13{]} that has passed the second reading in House of Lords will be
successfully implemented in the proposed date on the 1st of January 2027
{[}14{]}. The bill aims to prohibit the sale of tobacco to people born
on or after 1 January 2009 and provision about the licensing of retail
sales and the registration of retailers, alongside enforcement measures
and financial penalties for non-compliance. While some public debate has
characterised the policy as a form of ``nanny state'' intervention, the
Bill represents one of the most ambitious tobacco control measures in
recent UK history, aiming to protect future generations from nicotine
addiction and related harms.

Therefore, this analysis introduces the Tobacco and Vapes Bill within
the 2027 model cycle, assuming that individuals turning 18 in 2027 will
be the first cohort affected by the policy. The model applies a
year-on-year reduction in instigation rates, meaning the effect
compounds over time as younger cohorts progressively enter the
population under the new legislation. However, as there is currently no
empirical evidence quantifying the potential impact of the Bill, the
effect must be modelled through assumptions. Accordingly, Scenarios 1--4
represent varying levels of reduction in smoking instigation rates,
modelled as 10\%, 30\%, 60\%, and 90\% decreases, respectively, relative
to the baseline scenario of no policy intervention.

The results show that for those among 14 to 30 years old, smoking
prevalence is projected to decline to 4.13\% by 2047 under Scenario 1
(10\% reduction in instigation rate), while Scenario 2 (30\% reduction)
reaches 2.1\% in the same year. Scenarios 3 and 4, representing larger
reductions of 60\% and 90\%, produced almost identical trends, with
prevalence stabilising at 1.97\% by 2047.

{The latest results show that for those among 14 to 30 years old,
smoking prevalence is projected to decline to 3.77\% by 2047 under
Scenario 1 (10\% reduction in instigation rate), while Scenario 2 (30\%
reduction) reaches 1.91\% in the same year. Scenarios 3 and 4,
representing larger reductions of 60\% and 90\%, produced almost
identical trends, with prevalence stabilising at 1.78\% by 2047.}

\pandocbounded{\includegraphics[keepaspectratio]{smokingprevalence_reading_document_files/figure-pdf/unnamed-chunk-3-1.png}}

\pandocbounded{\includegraphics[keepaspectratio]{smokingprevalence_reading_document_files/figure-pdf/unnamed-chunk-4-1.png}}

The results show that for adults (18+), smoking prevalence is projected
to decline to 3.98\% by 2047 under Scenario 1 (10\% reduction in
instigation rate), while Scenario 2 (30\% reduction) reaches 2.97\% in
the same year. Scenarios 3 and 4, representing larger reductions of 60\%
and 90\%, produced almost identical trends, with prevalence stabilising
at approximately 2.7\% by 2047.

{The results for latest data show that for adults (18+), smoking
prevalence is projected to decline to 3.84\% by 2047 under Scenario 1
(10\% reduction in instigation rate), while Scenario 2 (30\% reduction)
reaches 2.93\% in the same year. Scenarios 3 and 4, representing larger
reductions of 60\% and 90\%, produced almost identical trends, with
prevalence stabilising at approximately 2.68\% by 2047.}

\pandocbounded{\includegraphics[keepaspectratio]{smokingprevalence_reading_document_files/figure-pdf/unnamed-chunk-5-1.png}}

\pandocbounded{\includegraphics[keepaspectratio]{smokingprevalence_reading_document_files/figure-pdf/unnamed-chunk-6-1.png}}

\section{Limitation}\label{limitation}

In developing the model, several assumptions and simplifications were
necessary, and as such, the findings should be interpreted in light of
these limitations.

\textbf{Firstly}, The model does not explicitly incorporate vaping
behaviour, if vaping became a substitute for smoking in the younger
cohorts, the model may overstate the decline in nicotine consumption,
since no information of vaping is incorporated to the model, individuals
who vape are classified as non-smokers within the current framework.
\textbf{Secondly}, the model does not incorporate information on how
long former smokers have been abstinent. All individuals entering the
ex-smoker state are assumed to have quit for one year at the start of
each cycle, which may overestimate their likelihood of relapse compared
to those who have remained smoke-free for longer periods. The model
applies England's average smoking prevalence to all net migrants, even
though individuals arriving from other countries may have distinct
smoking behaviours and cessation histories. Given this limitation, net
migration may be negative for ex-smoker states, particularly among older
age groups, reflecting Birmingham's overall pattern of net out-migration
at older ages. However, because negative values cannot be represented
within Markov state transitions, the model constrains these values to
zero. This adjustment ensures model stability but may slightly
underestimate the outflow of older former smokers in the projected
population structure. Fortunately, the extent of net out-migration of
ex-smoker states among older age groups was relatively modest, and
therefore is unlikely to have a substantial impact on the overall model
outcomes. \textbf{Thirdly}, parameter uncertainty for transition
probabilities can not be assessed, as the data sources used did not
provide associated confidence intervals or measures of variability.

\phantomsection\label{refs}
\begin{CSLReferences}{0}{0}
\bibitem[\citeproctext]{ref-ONS_2025_subnational_population_projections}
\CSLLeftMargin{{[}1{]} }%
\CSLRightInline{Office for National Statistics.
\href{https://www.ons.gov.uk/peoplepopulationandcommunity/populationandmigration/populationprojections/bulletins/subnationalpopulationprojectionsforengland/2022based}{Subnational
population projections for england: 2022‑based}. Office for National
Statistics, UK; 2025.}

\bibitem[\citeproctext]{ref-ONS_2023_adult_smoking_habits_gb}
\CSLLeftMargin{{[}2{]} }%
\CSLRightInline{Office for National Statistics.
\href{https://www.ons.gov.uk/peoplepopulationandcommunity/healthandsocialcare/healthandlifeexpectancies/bulletins/adultsmokinghabitsingreatbritain/2022}{Adult
smoking habits in the UK: 2022}. Office for National Statistics, UK;
2023.}

\bibitem[\citeproctext]{ref-NHS_Digital_2023_sdd_youth_england}
\CSLLeftMargin{{[}3{]} }%
\CSLRightInline{NHS Digital.
\href{https://digital.nhs.uk/data-and-information/publications/statistical/smoking-drinking-and-drug-use-among-young-people-in-england/2023}{Smoking,
drinking and drug use among young people in england, 2023}. NHS Digital;
2023.}

\bibitem[\citeproctext]{ref-ASH_2021_profile_16_17_smokers}
\CSLLeftMargin{{[}4{]} }%
\CSLRightInline{Action on Smoking and Health (ASH).
\href{https://ash.org.uk/uploads/Profile-of-16-and-17-year-old-smokers.pdf}{Profile
of 16/17‑year‑old smokers}. ASH; 2021.}

\bibitem[\citeproctext]{ref-ONS_2024_adult_smoking_habits_gb}
\CSLLeftMargin{{[}5{]} }%
\CSLRightInline{Office for National Statistics.
\href{https://www.ons.gov.uk/peoplepopulationandcommunity/healthandsocialcare/healthandlifeexpectancies/bulletins/adultsmokinghabitsingreatbritain/2024}{Adult
smoking habits in the UK: 2024}. Office for National Statistics, UK;
2025.}

\bibitem[\citeproctext]{ref-ONS_2024_national_life_tables}
\CSLLeftMargin{{[}6{]} }%
\CSLRightInline{Office for National Statistics.
\href{https://www.ons.gov.uk/peoplepopulationandcommunity/birthsdeathsandmarriages/lifeexpectancies/bulletins/nationallifetablesunitedkingdom/2021to2023additionaldata}{National
life tables, UK: 2021 to 2023 (additional data)}. Office for National
Statistics; 2024.}

\bibitem[\citeproctext]{ref-Song380}
\CSLLeftMargin{{[}7{]} }%
\CSLRightInline{Song F, Elwell-Sutton T, Naughton F, Gentry S. Future
smoking prevalence by socioeconomic status in england: A computational
modelling study. Tobacco Control 2021;30:380--5.
\url{https://doi.org/10.1136/tobaccocontrol-2019-055490}.}

\bibitem[\citeproctext]{ref-JEON2023S53}
\CSLLeftMargin{{[}8{]} }%
\CSLRightInline{Jeon J, Inoue-Choi M, Mok Y, McNeel TS, Tam J, Freedman
ND, et al. Mortality relative risks by smoking, race/ethnicity, and
education. American Journal of Preventive Medicine 2023;64:S53--62.
https://doi.org/\url{https://doi.org/10.1016/j.amepre.2022.12.006}.}

\bibitem[\citeproctext]{ref-Gillespie_et_al_smktrans}
\CSLLeftMargin{{[}9{]} }%
\CSLRightInline{Gillespie D, Webster L, Leeming G, Opazo-Breton M,
Brennan A.
\href{https://stapm.gitlab.io/smoking_state_transitions.html}{Smktrans:
An r package for estimating smoking state transition probabilities}
2023.}

\bibitem[\citeproctext]{ref-DHSC_2023_smokefree_generation_modeling}
\CSLLeftMargin{{[}10{]} }%
\CSLRightInline{Department of Health and Social Care.
\href{https://www.gov.uk/government/publications/smokefree-generation-policy-modelling-report/modelling-for-the-smokefree-generation-policy}{Modelling
for the smokefree generation policy}. UK Government; 2023.}

\bibitem[\citeproctext]{ref-CISNET}
\CSLLeftMargin{{[}11{]} }%
\CSLRightInline{Cancer Intervention and Surveillance Modeling Network
(CISNET).
\href{https://resources.cisnet.cancer.gov/projects/\#shg/iomr/resources}{Smoking
history generator (SHG) and individual outcome model resources} 2024.}

\bibitem[\citeproctext]{ref-CRUK_2021_smoking_prevalence_projections}
\CSLLeftMargin{{[}12{]} }%
\CSLRightInline{Cancer Research UK.
\href{https://www.cancerresearchuk.org/sites/default/files/cancer_research_uk_smoking_prevalence_projections_england_0.pdf}{Smoking
prevalence projections for england}. Cancer Research UK; 2021.}

\bibitem[\citeproctext]{ref-UK_Parliament_2024_Tobacco_Vapes_Bill}
\CSLLeftMargin{{[}13{]} }%
\CSLRightInline{UK Parliament.
\href{https://bills.parliament.uk/bills/3879}{Tobacco and vapes bill}
2024.}

\bibitem[\citeproctext]{ref-DHSC_2024_tobacco_vapes_policy}
\CSLLeftMargin{{[}14{]} }%
\CSLRightInline{Department of Health and Social Care.
\href{https://www.gov.uk/government/publications/the-tobacco-and-vapes-bill-creating-a-smoke-free-uk-and-tackling-youth-vaping/tobacco-and-vapes-bill-creating-a-smoke-free-uk-and-tackling-youth-vaping}{Tobacco
and vapes bill: Creating a smoke-free UK and tackling youth vaping}
2024.}

\end{CSLReferences}




\end{document}
